\documentclass[11pt,a4paper]{article}
\usepackage[utf8]{inputenc}
\usepackage[T1]{fontenc}
\usepackage[french]{babel}
\usepackage{lmodern,microtype}
\usepackage[margin=24mm]{geometry}
\usepackage{booktabs,tabularx,array,amsmath,enumitem}
\usepackage{hyperref}
\hypersetup{colorlinks=true,linkcolor=black,citecolor=black,urlcolor=black,
 pdftitle={Mesurer la valeur marginale des LLM face à des audits simples},
 pdfauthor={Étude LANCE, rédaction assistée par Codex}}
\setlength{\parindent}{0pt}
\setlength{\parskip}{4pt}
\setlength{\emergencystretch}{2em}
\clubpenalty=10000 \widowpenalty=10000
\setlist{nosep,leftmargin=*}
\setcounter{tocdepth}{1}
\newcolumntype{Y}{>{\raggedright\arraybackslash}X}
\newcommand{\fichier}[1]{\nolinkurl{#1}}
\newcommand{\revision}{2026-10-07 14:43 (Europe/Brussels, UTC+02:00)}
\title{Mesurer la valeur marginale des LLM\\face à des audits simples\\
 \large Diagnostic de D1 et six expériences hors ligne pour LANCE}
\author{Étude LANCE --- rédaction assistée par Codex}
\date{Dernière révision : \revision\\
 \small Statut : recherche exploratoire et protocoles proposés ; aucun résultat comparatif}
\begin{document}
\maketitle
\begin{abstract}
Le comité de thèse demande d'établir ce que des méthodes simples sans modèle
de langage accomplissent déjà, puis de mesurer le bénéfice de l'ajout du modèle.
Une lecture statique du dépôt montre un socle déterministe partagé et une
référence D1 conditionnelle, limitée à l'analyse et à la vérification. Les autres
comparateurs du protocole restent proposés ; les conflits de la copie locale
empêchent de présumer une exécution actuellement valide. Six expériences sur
observations et configurations fournies sont définies : routine, faible
prévalence, autorisation, information répartie, variations équivalentes et
restitution humaine. Elles séparent interprétation et collecte, comparent des
références simples crédibles, prévoient un oracle indépendant et conservent les
échecs. La décision peut être de garder les règles, d'utiliser le LLM seulement
pour une tâche délimitée, ou de conclure à une incertitude. Aucune marge métier,
taille d'échantillon ou performance n'est encore validée.
\end{abstract}
\tableofcontents

\section{Question, périmètre et méthode}
La question est : \textbf{pour quelles tâches le modèle apporte-t-il un gain
utile au-delà d'une méthode simple correctement conçue, à quel coût et avec
quelles erreurs ?} L'Internet des objets est abrégé IoT ; un modèle de langage
est abrégé LLM. Un audit peut réussir grâce à des outils et règles communs :
sa réussite absolue ne démontre pas une contribution du modèle.

Cette étude prolonge le protocole du 27 septembre~\cite{local} et répond à la
demande de comparateurs plus simples. Elle porte d'abord sur des dossiers
enregistrés ou des configurations documentées, sans nouvelle action sur des
systèmes. Les expériences sont des \emph{propositions}, non des résultats.
Elles ne mesurent ni découverte autonome, ni réussite d'intrusion, ni couverture
d'un réseau inconnu. Aucun audit de laboratoire ou appel au LLM applicatif
n'a été exécuté pour cette recherche.

La méthode associe lecture des index, des protocoles et des guides CLI/campagne,
inspection des modules cités au tableau~\ref{tab:code}, et recherche
méthodologique primaire ciblée. La tête Git consultée est
\fichier{a8dbeb9404f0cad960d61e42f02a1d70ca143a59}, avec modifications locales
et conflits : ce SHA ne suffit pas à identifier le contenu lu. Le fichier
\fichier{sans-llm-experiences-sources.json} conserve les empreintes des
modules inspectés et les périmètres de lecture des références externes, consultées
le 7 octobre 2026. Les résultats historiques de tests mentionnés dans les guides
n'ont pas été reproduits. Aucune importation applicative n'a été tentée.

Les cinq références méthodologiques concernent l'attribution causale,
l'indépendance des observations, l'équivalence, les oracles de test et la
maintenance. La lecture est ciblée : sections précisées aux références,
sans reproduction des articles. Le transfert de ces principes à LANCE est
\emph{notre proposition} ; les sources ne valident aucun gain de LANCE.

\section{Constats : ce que couvre réellement la référence D1}
\subsection{Un socle réel, mais un périmètre restreint}
D1 désigne l'automatisation conditionnelle sans LLM ; A1 désigne l'agent
adaptatif avec LLM. Le contrat expérimental disponible fixe un inventaire
public commun, le profil \texttt{full}, un worker et les phases 3--4 : analyse
puis vérification. La reconnaissance et la rédaction ne font pas partie de
ce contraste. En phase 3, les deux politiques utilisent d'abord le scanner
déterministe ; D1 évite ensuite l'analyse par agents LLM. Les sources visibles
étayent cette organisation, pas une qualité empirique des résultats.

\begin{table}[htbp]
\caption{Diagnostic statique de la copie consultée ; aucun taux de réussite n'est déduit.}
\label{tab:code}
\small
\begin{tabularx}{\linewidth}{@{}p{.23\linewidth}Y Y@{}}
\toprule
Objet et fichier & Fait visible & Limite pour la comparaison \\
\midrule
Entrée et périmètre\newline\fichier{audit_experiment.py}
& Valide l'inventaire public ; impose \texttt{full}, phases 3--4 et plafonds ; version D1 v2 (identifiant complet ci-dessous).
& Texte public et rôle peuvent révéler une solution ; l'empreinte d'environnement reste déclarative. \\
Phase 3\newline\fichier{analysis/scan.py}, \fichier{analysis/devices.py}
& Scanner et suivis déterministes communs ; les agents d'analyse sont évités en \texttt{rules}.
& L'apport d'A1 est un ajout au socle. Les observations finales d'A1 peuvent donc différer de celles de D1. \\
Extraction\newline\fichier{scanner.py}, \fichier{analysis/aggregation_loading.py}
& Extracteurs HTTP, MQTT, SSH, FTP, Redis, SNMP, CoAP, API/PKI/OTA/cloud/OT ; filtres de promotion selon profil et rôle.
& Présence d'un extracteur $\neq$ couverture démontrée. Certaines règles dépendent des simulateurs et de rôles connus. \\
Phase 4\newline\fichier{verification/rules.py}
& Plan préexistant ; une reprise admissible et une lecture HTTP complémentaire au plus ; trois appels au maximum par candidat dans cette boucle.
& Le plafond exclut les scans de phase 3. Un plan absent donne un cas non pris en charge, pas une conclusion saine. \\
Comparaison\newline\fichier{compare_policies.py}, \fichier{policy_campaign.py}
& Vérifie identité déclarée, artefacts et contrats ; conserve essais prévus et coûts observés, sépare paires admissibles.
& Pas d'intervalle par famille ; coût LLM partiel ; égalité des déclarations $\neq$ état réel identique. \\
\bottomrule
\end{tabularx}
\end{table}

La version déclarée des règles est
\fichier{shared-scanner-bounded-verification-v2}.
En phase 4, D1 reprend une opération répétable après une erreur transitoire
reconnue. Il peut remplacer une lecture HTTP insuffisante par une réponse
structurée sur la même URL. Il s'arrête sur confirmation, refus, erreur non
admissible ou preuve toujours insuffisante. Les verdicts utilisent le
synthétiseur de preuve commun. D1 est donc plus qu'un script linéaire,
mais ne représente pas toutes les méthodes classiques d'audit.

\subsection{Ce qui ne doit pas être annoncé comme acquis}
Des extracteurs couvrent déjà des autorisations API et des relations PKI.
Ce constat interdit de supposer que toute combinaison d'informations nécessite
un modèle. Il ne prouve pas que les cas nouveaux seront correctement traités.
Il faut distinguer l'existence d'une règle, son déclenchement, la production
d'un candidat, sa promotion et la preuve finale admise. La phase d'analyse
locale contient aussi un helper de compatibilité CVE, mais le chemin inspecté
de \fichier{analysis/scan.py} l'appelle sous condition \texttt{compact\_local}.
Cela ne démontre pas sa couverture dans l'expérience imposant \texttt{full} ;
il ne s'agit pas d'une preuve d'absence universelle de contrôle CVE.

Les fichiers \fichier{src/benchmark/metric_contract.py},
\fichier{src/agent/exploit_evidence.py} et
\fichier{src/agent/phases/verification/run.py} contiennent des marqueurs de
fusion. Les deux branches déclarent respectivement des contrats de mesure
\texttt{strict-v3.16}/\texttt{strict-v3.15} et de preuve
\texttt{evidence-v15}/\texttt{evidence-v14}. Aucun contrat unique n'est retenu
ici. Une analyse syntaxique statique sous Python 3.11.9 échoue aux marqueurs
de deux modules ; elle ne teste pas le runtime cible Python 3.12. Les
validations historiques ne garantissent donc pas l'importabilité actuelle.
La remise en état et les tests ciblés constituent une précondition d'une
future expérience intégrée, pas un travail réalisé dans cette recherche.

D0 (procédure fixe), A0 (plan LLM initial ensuite figé) et D1-R (D1 avec
rédaction par LLM après clôture technique) sont documentés comme propositions.
Le manifeste de 24 bras est une préparation, sans preuve de campagne réalisée.
Les six expériences suivantes n'utilisent pas ces nombres comme effectif validé.

\section{Comparateurs simples et attribution de l'effet}
\subsection{Construire une référence crédible avant d'ajouter de la complexité}
Le tableau~\ref{tab:arms} conserve les noms du protocole. Pour les expériences
hors ligne, « interpréteur à règles » et « interpréteur LLM » sont des lecteurs
de dossiers, pas des nouveaux noms pour D1 et A1. L'étude de méthodes d'audit
du 7 octobre utilise séparément E0, E1 et E2 pour des représentations de preuves ;
ces identifiants ne sont ni renommés, ni assimilés à D0/D1/A0/A1.

\begin{table}[htbp]
\caption{Ordre proposé des comparateurs ; les lecteurs hors ligne restent distincts des contrôleurs.}
\label{tab:arms}
\small
\begin{tabularx}{\linewidth}{@{}p{.13\linewidth}Y Y@{}}
\toprule
Référence & Définition intelligible & Question et priorité \\
\midrule
D0 & Procédure fixe, sortie structurée et règles de constat prédéfinies. & Établir ce qui fonctionne sans adaptation ; référence secondaire. \\
D1 & Règles conditionnelles, état observable, reprises bornées et arrêt explicite. & Référence principale. Renforcer les lacunes connues sur développement, sans voir l'oracle du test. \\
Lecteurs & Même dossier normalisé, règles contre LLM, sans outil supplémentaire. & Premier contraste pour l'interprétation ; décrit séparément du pipeline interactif. \\
D1-R & Mêmes résultats techniques D1, rapport rédigé par LLM après gel. & Expérience de restitution ; aucun gain de détection attribuable au rapport. \\
A0/A1 & Plan LLM figé contre agent adaptatif. & Contraste ultérieur si le besoin de décision interactive est établi ; non testé ici. \\
\bottomrule
\end{tabularx}
\end{table}

La priorité est une table de règles et une normalisation commune, avec une
matrice de couverture lisible : propriété, données requises, cas pris en charge,
motif d'abstention et tests positifs/négatifs. Des jointures relationnelles
simples ou une politique d'accès explicite sont ajoutées seulement si les
cas exigent ces informations. Une règle dépendant de l'identifiant privé d'un
scénario est inadmissible. Une règle métier publique accordée au LLM doit
également être disponible au comparateur simple.

Avant le gel, un relecteur différent des auteurs examine les lacunes réparables
de D1 et les connaissances liées aux simulateurs. Archiver les heures de
développement des règles, des parseurs, des prompts et de leurs essais. Le
test final reste inaccessible à ces auteurs. Une comparaison avec un outil
externe mesure ensuite l'utilité de systèmes complets ; des différences de
couverture ou de collecte empêchent d'attribuer tout l'écart au LLM.

\subsection{Trois effets à ne pas mélanger}
La première question est l'interprétation : même information, même format
accessible, même propriété attendue et même budget d'analyse ; seule la méthode
de lecture change. La seconde concerne le choix d'observations dans une future
expérience intégrée. La troisième est le travail humain de restitution.
La distinction entre effet d'une intervention définie et association motive
ces contrastes ; notre application au logiciel est une adaptation méthodologique
de la discussion sur la randomisation de Hernán et Robins~\cite{hernan}.

Un dossier construit uniquement à partir des candidats de LANCE manque les
cas que le pipeline n'a jamais proposés. L'échantillonnage doit partir des
propriétés et configurations indépendantes, puis inclure cas détectés,
non détectés, sains et indéterminés. Quand des observations proviennent d'un
run A1 historique, préciser leur provenance : une comparaison de lecteurs sur
ces mêmes observations ne mesure plus l'avantage de leur collecte.

\section{Protocole commun aux six expériences proposées}
Les dossiers contiennent des observations expurgées de secrets réels et un
contexte public documenté : rôle, ressource, politique, date et provenance.
L'oracle privé conserve \textbf{deux étiquettes distinctes}. La première,
$T_{\text{état}}$, décrit la propriété dans l'état de référence indépendant :
« vrai », « faux » ou « inconnu ». La seconde, $T_{\text{preuve}}$, décrit
ce que les observations du dossier permettent d'établir : « confirmable »,
« réfutable » ou « insuffisant ». Une propriété vraie dans la configuration
peut rester insuffisamment documentée ; une abstention est alors correcte
pour le lecteur. Un conflit non résolu peut aussi empêcher une conclusion.
Une information manquante n'efface jamais l'étiquette d'état connue.

Barr et al. distinguent oracle partiel, vérité terrain et limites de justesse
et de complétude~\cite{barr}. Pour LANCE, proposer un double contrôle :
conformité du verdict à la propriété réellement configurée et adéquation de
ses références aux observations disponibles. Deux annotateurs, distincts
des auteurs des méthodes, travaillent avec identifiants de méthode masqués ;
les désaccords sont conservés puis arbitrés. La consigne et les cas litigieux
sont versionnés. Le validateur commun du pipeline est un facteur de symétrie,
pas l'unique oracle indépendant.

Chaque sortie comprend identifiant de propriété, verdict, références de
preuves, contradictions, informations manquantes et motif d'abstention.
Les références doivent pointer vers des observations fournies, pertinentes
pour cette ressource, à cette date. Aucun score de confiance auto-déclaré
ne remplace ce contrôle. Les métriques hors ligne sont nouvelles et doivent
être documentées ; elles ne sont pas présentées comme le contrat courant
\texttt{strict-v3}.

Le rappel d'état utilise comme dénominateur tous les positifs connus
indépendamment sélectionnés, même si leur dossier est incomplet. Il décrit la
restitution de cette population, affectée par la couverture des observations.
Le rappel conditionnel du lecteur utilise seulement les positifs connus dont
la preuve est confirmable depuis le dossier. Une confirmation créditée exige
un verdict juste et une référence de preuve adéquate. Publier les deux
dénominateurs, la couverture documentaire,
les abstentions et les erreurs : une abstention correcte sur preuve insuffisante
reste visible dans le bilan global, sans être imputée à une faute d'interprétation.
Les états inconnus et dossiers incomplets sont conservés avec leur statut.

L'unité de décision est la propriété d'un dossier ; l'unité de généralisation
est sa famille d'origine, définie par générateur, services et structures
techniques partagés, plutôt que par identifiant de scénario. Les répétitions
LLM sont imbriquées dans le dossier, les variations dans la famille. Fixer
répartition développement/test par famille avant de concevoir règles et prompts.
Un test réservé aux auteurs ne certifie pas l'absence de contamination du
préentraînement lorsque le corpus d'origine est public. Le problème de pseudoréplication
décrit par Lazic motive cette prudence~\cite{lazic} ; ses résultats biologiques
ne sont pas des estimations de dépendance pour notre corpus. Une nouvelle
graine ou copie JSON ne crée pas une famille indépendante.

Les bras reçoivent exactement les mêmes octets utiles ou une normalisation
équivalente vérifiée ; ordre d'exécution et de présentation sont équilibrés
ou randomisés. Journaliser les coûts de toutes les tentatives, les erreurs de
format, arrêts, exclusions et valeurs inconnues. L'absence de réseau dans une
expérience hors ligne ne prouve aucune disponibilité des outils réels.
Un budget total de temps/calcul est comparable ; les tokens servent uniquement
à comptabiliser le bras qui utilise un modèle.

\subsection{Fiche 1 --- Contrôles routiniers}
\textbf{Question.} Une procédure fixe ou des règles conditionnelles suffisent-elles
à distinguer une configuration conforme d'une non-conformité connue ?
\textbf{Données et entrée.} Dossiers de développement documentant configurations
de service et réponses déjà enregistrées ; propriétés stables, cas conformes,
non conformes et réponses incomplètes, y compris cas absents des candidats LANCE.
\textbf{Sortie et oracle.} Verdict et preuve minimale pour chaque propriété,
contrôlés contre configuration et observations indépendamment annotées.

\textbf{Comparateurs et contrôles.} Lecteur à règles fixes, lecteur conditionnel,
puis lecteur LLM sur les mêmes dossiers ; sortie et normalisation communes.
Un service simplement joignable ne constitue pas à lui seul une non-conformité.
Les connaissances métier sont publiques et identiques.
\textbf{Mesures et coûts.} Erreurs par propriété, précision et rappels d'état
et conditionnel avec les dénominateurs définis ci-dessus, abstentions correctes,
taux de preuves attribuées et temps/calcul,
avec coût de construction des règles.
\textbf{Arrêt et poursuite.} Suspendre le test si une différence de format ou
de couverture explique le contraste. Poursuivre vers test réservé si les
méthodes simples couvrent correctement les cas de développement ; rechercher
alors leur éventuelle non-infériorité plutôt qu'une victoire imposée du LLM.

\subsection{Fiche 2 --- Contrôles sains et faible prévalence}
\textbf{Question.} Le LLM apporte-t-il des alertes utiles sans produire davantage
de fausses alertes lorsque les non-conformités sont rares ?
\textbf{Données et entrée.} Dossiers conformes validés et proches cas non conformes,
à partir d'une population définie indépendamment des sorties des méthodes.
Prévalences par strate fixées avant test ; aucune distribution synthétique
ne représente automatiquement un déploiement réel.
\textbf{Sortie et oracle.} Liste d'alertes et abstentions avec preuve ; oracle
privé vérifiant l'état conforme et les contrôles effectivement observés.

\textbf{Comparateurs et contrôles.} Lecteurs fixes, conditionnels et LLM ; même
information. Ajouter réponses d'authentification refusée, contenu public
inoffensif et observation manquante. Un dossier non analysé n'est pas un sain
correctement reconnu.
\textbf{Mesures et coûts.} Fausses alertes par dossier et par propriété,
spécificité avec dénominateur explicite, valeur prédictive positive à la
prévalence déclarée, rappels d'état et conditionnel des rares cas positifs,
et temps de tri humain. Les fausses alertes et la spécificité sont les mesures
principales des contrôles sains. Pour $F1=2TP/(2TP+FP+FN)$, une strate saine
avec $FP>0$ donne zéro ; avec dénominateur nul, annoncer la convention choisie
(valeur indéfinie dans ce protocole). Le silence universel est pénalisé sur
les cas positifs.
\textbf{Arrêt et poursuite.} Suspendre si les contrôles sains ne sont pas validés
ou si l'extraction des cas dépend des alertes LANCE. Étendre seulement si la
marge de fausses alertes préfixée peut être évaluée avec une précision utile.

\subsection{Fiche 3 --- Autorisation contextuelle}
\textbf{Question.} Une politique explicite permet-elle des conclusions aussi
justes qu'une lecture LLM du contexte d'autorisation ?
\textbf{Données et entrée.} Matrice publique sujet--ressource--action--tenant,
configuration fournie et traces préexistantes de décisions d'accès. Des cas
partagent la même réponse technique mais des politiques légitimes différentes.
\textbf{Sortie et oracle.} Relation entre accès observé et accès attendu,
preuve associant identité, ressource et politique ; statut indéterminé si
l'identité ou la portée de la politique n'est pas connue.

\textbf{Comparateurs et contrôles.} Règles conditionnelles avec table de politique,
contre lecteur LLM recevant cette même table. D0 sert de repère, mais n'est
pas le seul comparateur. Cas autorisé, refus conforme, accès contraire à la
politique et identité absente sont présents ; les réponses HTTP seules ne
déterminent pas le droit d'accès.
\textbf{Mesures et coûts.} Classification de conformité, erreurs de liaison
identité/ressource, abstention et coût d'encodage de la politique.
\textbf{Arrêt et poursuite.} Suspendre si l'entrée publique contient l'étiquette
privée de non-conformité. La matrice normative de droits est une entrée légitime
et nécessaire aux deux bras ; sa présence n'est pas une fuite d'oracle.
Généraliser uniquement à de nouvelles politiques réservées ; une victoire
contre des règles privées du contexte ne serait pas un effet du modèle isolé.

\subsection{Fiche 4 --- Information répartie et contradictions}
\textbf{Question.} Le LLM relie-t-il mieux des observations et repère-t-il des
contradictions au-delà de jointures simples ?
\textbf{Données et entrée.} Dossiers avec plusieurs observations référencées,
identifiants de ressource et temps ; cas cohérents, contradictoires, périmés
et incomplets. L'oracle associe les observations à la propriété pertinente,
sans se limiter aux conclusions déjà proposées par LANCE.
\textbf{Sortie et oracle.} Verdict, liens de preuve, contradictions et manque
d'information ; l'oracle distingue conflit réel, ressources différentes et
changement d'état documenté.

\textbf{Comparateurs et contrôles.} Règles conditionnelles sur un index tabulaire,
puis, si besoin, jointures déterministes explicites ; lecteur LLM avec mêmes
observations, clés, obligations et budget de lecture. Si un graphe est étudié,
le contraste de représentation E0/E1 reste séparé de celui du lecteur.
\textbf{Mesures et coûts.} Attribution valide, fausse fusion, contradictions
repérées ou inventées, abstentions justifiées et temps/calcul ; coût de maintien
des clés et des relations.
\textbf{Arrêt et poursuite.} Suspendre si une méthode possède des obligations
ou des observations supplémentaires. Garder l'index simple si les relations
n'apportent aucun gain utile ; tester le lecteur LLM sur des compositions
réservées seulement après une référence relationnelle crédible.

\subsection{Fiche 5 --- Variations équivalentes}
\textbf{Question.} La qualité résiste-t-elle à une présentation différente
qui préserve les mêmes faits et la même propriété de sécurité ?
\textbf{Données et entrée.} Paires original/variation, approuvées indépendamment :
ordre des champs, libellés publics équivalents, répartition en fichiers ou
formulations de politique. Préserver identifiants et références ou publier
leur correspondance ; ne pas retirer une information à un seul bras.
\textbf{Sortie et oracle.} Même verdict attendu et preuves correspondantes,
avec conservation explicite des cas indéterminés.

\textbf{Comparateurs et contrôles.} Lecteurs à règles et LLM sur observations
structurées communes ; seconde strate brute seulement pour tester les parseurs.
Le test métamorphique exprime une relation attendue entre deux exécutions~\cite{barr}.
Il ne prouve pas qu'un verdict invariant est juste : les deux versions peuvent
partager la même erreur.
\textbf{Mesures et coûts.} Justesse absolue de chaque version, taux d'invariance,
pertes de preuve et variation de l'écart LLM--règles par famille ; coûts de
normalisation et de maintenance.
\textbf{Arrêt et poursuite.} Exclure avec motif toute transformation ne préservant
pas les faits. Dans la strate normalisée, une perte d'information due au parseur
invalide le contraste sémantique ; suspendre ce contraste et corriger sa préparation.
Dans la strate brute, l'échec de parsing reste un résultat de robustesse et
son coût est conservé. Une réussite autorise une conclusion de robustesse aux
transformations testées, pas de généralisation à de nouveaux protocoles.

\subsection{Fiche 6 --- Rapport et temps humain}
\textbf{Question.} D1-R réduit-il le travail d'un auditeur à résultats techniques
constants, sans augmenter les erreurs ?
\textbf{Données et entrée.} Dossiers techniques gelés, comprenant preuves,
limites et conclusions ; rapports par gabarit D1 et rédaction LLM D1-R.
Tâches d'auditeur : retrouver une preuve, corriger un constat, prioriser une
remédiation documentée. Plusieurs dossiers de familles différentes sont utilisés.
\textbf{Sortie et oracle.} Réponse finale de l'auditeur et corrections, évaluées
contre les mêmes résultats techniques ; aucune modification des détections.

\textbf{Comparateurs et contrôles.} Dossiers sources identiques et contraintes
de production communes, notamment de longueur ; ordre équilibré, expérience des auditeurs consignée,
affectation évitant qu'un lecteur revoie le même dossier avec l'autre rapport.
Masquer l'origine des rapports autant que possible sans prétendre un aveugle
parfait. Les effets auditeur et dossier restent distincts.
\textbf{Mesures et coûts.} Temps actif, réussite des tâches, corrections, omission
de réserves et affirmations inventées ; latence, consommation du modèle et
temps de préparation/revue. Satisfaction seule ne suffit pas.
\textbf{Arrêt et poursuite.} Suspendre si les dossiers sources ou les contraintes
de production diffèrent, ou si le recrutement ne permet qu'un avis de l'auteur.
Les omissions et affirmations inventées dans les rapports sont des erreurs
mesurées et conservées dans le bilan, avec leurs coûts de revue et correction ;
elles ne justifient pas l'exclusion d'un rapport. Cette tâche de lecture
n'entraîne aucune décision sur un système réel. Conserver D1-R seulement si
un gain humain utile est compatible avec les marges d'erreur ; le résultat
concerne la restitution, indépendamment de la détection.

\section{Mesures, décision et séquence réaliste}
\subsection{Un effet utile, une équivalence ou une conclusion indéterminée}
Fixer une métrique principale et une portée par expérience avant le test final.
Publier aussi les nombres bruts, les contrôles sains, les erreurs de preuve,
les abstentions et la fraction de dossiers évaluables. Choisir les marges à
partir du besoin d'audit et du développement : gain minimal utile,
dégradation acceptable, coût et temps humain. Ce rapport ne leur attribue
aucune valeur numérique prétendument validée.

Pour une métrique où une valeur élevée est meilleure, poser
$\Delta=Q_{\mathrm{LLM}}-Q_{\text{règles}}$. Si la borne inférieure de
l'intervalle dépasse la marge utile préfixée, le gain peut soutenir l'usage du
modèle, sous contraintes d'erreurs et de coût. Pour soutenir la non-infériorité des règles, l'hypothèse
est que $Q_{\text{règles}}-Q_{\mathrm{LLM}}$ reste au-dessus de la perte
maximale tolérée négative. L'équivalence est bilatérale, autour de marges fixées
avant test : l'intervalle approprié doit être entièrement compris entre ces
marges. Une différence non significative ne démontre aucune de ces deux
propriétés~\cite{lakens}. Le choix de l'estimateur et de l'intervalle doit
respecter les familles et le type de mesure ; copier un test $t$ générique
sur tous les constats dépendants serait insuffisant.

Proposer une analyse appariée, avec agrégation répétition $\rightarrow$ dossier
$\rightarrow$ famille et pondération des familles annoncée à l'avance.
Si les familles sont trop peu nombreuses pour l'inférence visée, publier un
bilan descriptif et une conclusion indéterminée. La taille du test devra être
justifiée après le développement, à partir de la variabilité observée et de
marges opérationnelles, avec analyse de précision ou simulation adaptée.
Le manifeste historique n'est pas une analyse de puissance.

Conserver toutes les tentatives dans le bilan opérationnel : échec de format,
arrêt, budget dépassé ou dossier incomplet. Les métriques techniques manquantes
restent inconnues, avec analyse de sensibilité prévue ; elles ne deviennent
pas silencieusement zéro. Les conclusions tirées seulement des paires
exploitables sont étiquetées conditionnelles. Les analyses secondaires et
leurs éventuelles corrections de multiplicité sont déclarées avant le test.

\subsection{Coût de construction, d'exploitation et de maintenance}
Le comparateur actuel annonce un coût d'usage du LLM, sans mesurer toute
l'infrastructure et le travail humain. Proposer le coût total explicite :
\[
 C(N,H)=C_{\text{construction}}+C_{\text{maintenance}}(H)
       +\sum_{i=1}^{N}(C_{\text{calcul},i}+C_{\text{modèle},i}
       +C_{\text{préparation/revue},i}).
\]
$N$ est le volume de dossiers ou d'audits et $H$ l'horizon déclaré. Les règles
et les prompts ont tous deux un coût de maintenance ; un prix API nul ne
signifie pas un coût nul. Présenter heures et montants séparément, avec taux
et scénarios d'amortissement documentés. Sculley et al. motivent le suivi des
dépendances, configurations et changements externes~\cite{sculley}, sans
fournir un tarif ou une supériorité économique applicable à LANCE.

\subsection{Ordre proposé et décisions possibles}
La première livraison expérimentale serait le corpus annoté et un lecteur
simple sur les fiches 1--2, puis le contraste de lecteurs sur observations
identiques. Les fiches 3--4 ajoutent le contexte seulement quand il correspond
à un besoin réel. La fiche 5 accompagne les comparaisons ; la fiche 6 nécessite
une étude humaine distincte. Après gel, évaluer sur familles réservées sans
retoucher règles et prompts sur les échecs du test. Une correction crée une
nouvelle version et nécessite un nouveau test réservé.

Trois décisions sont recevables : garder les règles sur les tâches où elles
suffisent selon les marges retenues ; réserver le LLM à une tâche pour laquelle
un gain utile est mesuré ; ou conserver le résultat « indéterminé » quand la
preuve manque. Si le modèle n'aide que la rédaction, D1-R peut être utile
sans bénéfice de décision technique. Si l'aide tient surtout au développement
de nouvelles règles, une étude séparée humain seul/humain assisté devra mesurer
ce bénéfice de construction, avec exécution finale déterministe.

\section{Limites et conclusion}
Cette recherche ne démontre aucune performance comparative. La lecture statique
ne prouve ni importabilité actuelle, ni disponibilité des outils, ni validité
des configurations réellement déployées. Les conflits sont conservés ; les
sources externes et les archives historiques ne réparent pas cette limite.
La comparaison de lecteurs ne permet pas d'inférer la valeur d'un contrôleur
adaptatif complet. Les connaissances préentraînées du LLM restent un facteur
inégalizable : le contraste mesure une méthode avec ses connaissances,
pas le raisonnement pur.

Les premiers tests doivent établir la force et les limites des méthodes
simples sur des dossiers indépendamment sélectionnés, avant de chercher des
cas favorables au modèle. Les six fiches rendent ce programme réfutable :
besoin défini, même information, oracle indépendant, coût complet et règle de
décision avant les résultats. Une conclusion négative ou conditionnelle peut
donner une contribution de thèse : délimiter les tâches où une automatisation
simple suffit et celles où un LLM apporte une valeur mesurée.

\begin{thebibliography}{9}
\interlinepenalty=10000
\addcontentsline{toc}{section}{Références}
\bibitem{local} LANCE. \emph{Agent / automatisation}, protocole du 27 septembre
2026, suivi des 27--28 septembre, guides CLI/campagne et modules cités.
Lecture statique le 7 octobre 2026 ; tests historiques non reproduits.
\fichier{research/2026-09-27-agent-vs-automation/protocol.md},
\fichier{docs/benchmark/agent-vs-automation-cli.md} et
\fichier{docs/benchmark/agent-vs-automation-campaign.md}.

\bibitem{hernan} M. A. Hernán, J. M. Robins. \emph{Causal Inference: What If}.
2020 ; version PDF du 26 avril 2024. Section 2.1 « Randomization » : extrait
indexé du PDF primaire consulté le 7 octobre 2026 ; fichier complet inaccessible
dans l'outil, autres chapitres non exploités.
\href{https://www.hsph.harvard.edu/miguel-hernan/wp-content/uploads/sites/1268/2024/04/hernanrobins_WhatIf_26apr24.pdf}{PDF primaire Harvard}.

\bibitem{lazic} S. E. Lazic. \emph{The problem of pseudoreplication in
neuroscientific studies: is it affecting your analysis?} BMC Neuroscience 11:5,
14 janvier 2010, doi:10.1186/1471-2202-11-5. Résumé, Background et sections
sur observations répétées/hiérarchiques lus le 7 octobre 2026 ; résultats de
neuroscience non transférés quantitativement.
\href{https://link.springer.com/article/10.1186/1471-2202-11-5}{Article éditeur}.

\bibitem{lakens} D. Lakens. \emph{Equivalence Tests: A Practical Primer for
t Tests, Correlations, and Meta-Analyses}. Social Psychological and Personality
Science 8(4):355--362, 2017, doi:10.1177/1948550617697177. Version éditeur
TU/e : Testing for Equivalence, Setting Equivalence Bounds et Discussion lus
le 7 octobre 2026 ; aucune procédure statistique exécutée, adaptation aux
mesures dépendantes de LANCE encore nécessaire.
\href{https://pure.tue.nl/ws/portalfiles/portal/80918653/lakeequi2017.pdf}{PDF éditeur déposé à TU/e}.

\bibitem{barr} E. T. Barr, M. Harman, P. McMinn, M. Shahbaz, S. Yoo.
\emph{The Oracle Problem in Software Testing: A Survey}. IEEE TSE
41(5):507--525, mai 2015, doi:10.1109/TSE.2014.2372785. Introduction,
définitions, soundness/completeness lus dans le PDF UCL le 7 octobre 2026 ;
techniques du corpus non reproduites. La métamorphie est présentée comme
oracle partiel, sans validation d'une implémentation LANCE.
\href{https://discovery.ucl.ac.uk/1471263/1/06963470.pdf}{PDF éditeur déposé à UCL}.

\bibitem{sculley} D. Sculley et al. \emph{Hidden Technical Debt in Machine
Learning Systems}. NIPS 28, 2015. Sections 3, 5--7 et 9 lues le 7 octobre 2026 ;
discussion d'ingénierie sans mesure de coût spécifique LANCE.
\href{https://proceedings.nips.cc/paper/2015/file/86df7dcfd896fcaf2674f757a2463eba-Paper.pdf}{PDF des actes}.
\end{thebibliography}
\end{document}
